% ══════════════════════════════════════════════════════════════
%  Chapter 13: Node Orchestration --- Coordinating the Runtime
% ══════════════════════════════════════════════════════════════
\chapter{Node Orchestration}
\label{ch:node-orchestration}

\begin{learnbox}
\begin{itemize}
  \item Coordinate blockchain state, mempool, network, mining, and validation through the \code{NodeContext} API
  \item Understand how the node module provides a unified interface over multiple subsystems
  \item Trace how an incoming request flows through the node context to the appropriate handler
  \item Explain why centralized orchestration simplifies the interaction between blockchain components
\end{itemize}
\end{learnbox}

\lettrine{A}{t} \index{runtime}runtime, several \index{subsystem}subsystems must work together:

\begin{itemize}
  \item \textbf{\index{network!layer}Network} receives bytes and produces typed \index{message}messages (\code{\index{Package}Package})
  \item \textbf{\index{NodeContext}NodeContext} decides what \index{subsystem}subsystem should handle the \index{message}message (\index{mempool}mempool, \index{chainstate}chainstate, \index{mining}mining)
  \item \textbf{\index{chainstate}Chainstate} persists \index{block}blocks and updates the \index{canonical}canonical \index{tip}tip + \index{UTXO}UTXO
  \item \textbf{\index{mempool}Mempool} stores \index{unconfirmed transaction}unconfirmed \index{transaction}transactions and marks \index{output}outputs as ``in \index{mempool}mempool''
  \item \textbf{\index{miner}Mining} turns \index{mempool}mempool contents into a \index{block}block and relays the new \index{tip}tip
  \item \textbf{\index{peer}Peers} is the \index{node}node's \index{peer}peer set used for relay
\end{itemize}

This chapter explains the \code{\index{Rust!module}bitcoin/src/node} \index{Rust!module}module as a \index{Rust!developer}Rust implementer reads it: as a \index{coordination}coordination layer that turns ``\index{network!message}network \index{message}messages and API calls'' into ``\index{chainstate}chainstate updates, \index{mempool}mempool updates, \index{mining}mining, and \index{peer!relay}peer relay''.

% ──────────────────────────────────────────────────────────────
\section{What ``Node Orchestration'' Means}
\label{sec:ch13-what-means}

The \index{runtime}runtime flow is:

\begin{minted}[fontsize=\footnotesize,linenos=false]{text}
TCP -> Package -> NodeContext -> (mempool | add\_block | mining | relay)
\end{minted}

The key \index{entry point}entry points are:

\begin{itemize}
  \item \code{\index{NodeContext!process\_transaction}NodeContext::process\_transaction} (\index{mempool}mempool admission)
  \item \code{\index{NodeContext!add\_block}NodeContext::add\_block} (\index{chain}chain extension from a \index{peer}peer)
  \item \code{\index{NodeContext!mine\_empty\_block}NodeContext::mine\_empty\_block} (local \index{mining}mining trigger)
\end{itemize}

These delegate to:

\begin{itemize}
  \item \code{\index{txmempool}txmempool::\{add\_to\_memory\_pool, remove\_from\_memory\_pool\}} for \index{mempool}mempool \index{state}state
  \item \code{\index{miner}miner::\{should\_trigger\_mining, process\_mine\_block, broadcast\_new\_block\}} for the \index{mining}mining \index{pipeline}pipeline
\end{itemize}

\begin{warnbox}[Critical Role]
\code{\index{NodeContext}NodeContext} is the single \index{entry point}entry point for all \index{blockchain}blockchain operations. Every user-facing \index{interface}interface --- the \index{web API}web API, the desktop \index{admin dashboard}admin UIs, and the \index{wallet}wallet applications --- communicates with the \index{blockchain}blockchain exclusively through \code{\index{NodeContext}NodeContext}. Understanding this \index{Rust!module}module is essential for understanding how the \index{system}system's components fit together.
\end{warnbox}

% ──────────────────────────────────────────────────────────────
\section{The Façade: \texttt{NodeContext} is the Coordinator}
\label{sec:ch13-facade}

\code{NodeContext} is the central façade used by other layers (network server, web/API) to interact with the node. It owns a \code{BlockchainService} handle and reaches out to:

\begin{itemize}
  \item the global mempool (\code{GLOBAL\_MEMORY\_POOL}) via \code{txmempool}
  \item the peer set (\code{GLOBAL\_NODES}) via \code{peers}
  \item the network send primitives (\code{send\_inv}) to announce inventory
  \item the mining pipeline (\code{miner}) to build blocks and broadcast new tips
\end{itemize}

% ──────────────────────────────────────────────────────────────
\subsection{Transaction Ingestion: Accept $\to$ Relay $\to$ Maybe Mine}
\label{sec:ch13-transaction-ingestion}

This is the path you will trace most often when debugging ``why didn't my transaction propagate / get mined?''

The orchestration call graph for transactions:

\begin{minted}[fontsize=\footnotesize,linenos=false]{text}
Inbound tx bytes
  -> (Network) Package::Tx
  -> NodeContext::process_transaction
      -> txmempool::add_to_memory_pool
      -> spawn background task:
         NodeContext::submit_transaction_for_mining
             -> (if central) broadcast INV(txid) to peers
             -> (if mining threshold) miner::process_mine_block
\end{minted}

\begin{listing}[htbp]
\caption{NodeContext initialization and transaction acceptance}
\label{lst:ch13-nodecontext-init}
\begin{minted}[fontsize=\footnotesize]{rust}
use crate::chain::{BlockchainService, UTXOSet};
use crate::error::{BtcError, Result};
use crate::node::miner::*;
use crate::node::txmempool::*;
use crate::node::GLOBAL_NODES;
use crate::{Block, Transaction, WalletAddress};
use std::net::SocketAddr;

#[derive(Clone, Debug)]
pub struct NodeContext {
    blockchain: BlockchainService,  // Only owned dependency
}

impl NodeContext {
    pub fn new(blockchain: BlockchainService) -> Self {
        Self { blockchain }
    }

    pub async fn btc_transaction(
        &self,
        wlt_frm_addr: &WalletAddress,
        wlt_to_addr: &WalletAddress,
        amount: i32,
    ) -> Result<String> {
        let utxo_set = UTXOSet::new(
            self.blockchain.clone());
        let utxo = Transaction::new_utxo_transaction(
            wlt_frm_addr,
            wlt_to_addr,
            amount,
            &utxo_set
        ).await?;
        let addr_from =
            crate::GLOBAL_CONFIG.get_node_addr();
        self.process_transaction(&addr_from, utxo)
            .await
    }

    pub async fn process_transaction(
        &self,
        addr_from: &SocketAddr,
        utxo: Transaction,
    ) -> Result<String> {
        // 1) Dedupe check
        if transaction_exists_in_pool(&utxo) {
            return Err(BtcError
                ::TransactionAlreadyExistsInMemoryPool(
                    utxo.get_tx_id_hex(),
                ));
        }

        // 2) Accept into mempool
        // (flags outputs as "in mempool")
        add_to_memory_pool(utxo.clone(), &self.blockchain)
            .await?;

        // 3) Fire-and-forget: broadcast + mining
        // in background task
        let context = self.clone();
        let addr_copy = *addr_from;
        let tx = utxo.clone();
        tokio::spawn(async move {
            let _ = context
                .submit_transaction_for_mining(
                    &addr_copy,
                    tx
                ).await;
        });

        Ok(utxo.get_tx_id_hex())
    }
}
\end{minted}
\end{listing}

After accepting a transaction into the mempool, the node spawns background tasks for relay and mining decisions:

\begin{listing}[htbp]
\caption{Background relay and mining submission}
\label{lst:ch13-submit-mining}
\begin{minted}[fontsize=\footnotesize]{rust}
impl NodeContext {
    async fn submit_transaction_for_mining(
        &self,
        addr_from: &SocketAddr,
        utxo: Transaction,
    ) -> Result<()> {
        // Background: relay to peers (if central node)
        // and maybe trigger mining
        let my_node_addr =
            GLOBAL_CONFIG.get_node_addr();
        if my_node_addr == crate::node::CENTERAL_NODE {
            let nodes = self
                .get_nodes_excluding_sender(addr_from)
                .await?;
            for node in &nodes {
                send_inv(
                    &node.get_addr(),
                    OpType::Tx,
                    &[utxo.get_id_bytes()],
                )
                .await;
            }
        }

        // Mining: trigger if mempool reaches threshold
        if should_trigger_mining() {
            if let Some(addr) =
                GLOBAL_CONFIG.get_mining_addr() {
                if let Ok(txs) =
                    prepare_mining_utxo(&addr) {
                    if !txs.is_empty() {
                        process_mine_block(
                            txs,
                            &self.blockchain
                        ).await?;
                    }
                }
            }
        }
        Ok(())
    }

    async fn get_nodes_excluding_sender(
        &self,
        addr_from: &SocketAddr,
    ) -> Result<Vec<Node>> {
        let nodes = GLOBAL_NODES
            .get_nodes()?
            .into_iter()
            .filter(|node| {
                let node_addr = node.get_addr();
                node_addr != *addr_from
                    && node_addr !=
                    GLOBAL_CONFIG.get_node_addr()
            })
            .collect();
        Ok(nodes)
    }
}
\end{minted}
\end{listing}

% ──────────────────────────────────────────────────────────────
\subsection{Block Ingestion and Chain/Peer Queries}
\label{sec:ch13-block-ingestion}

This is the subset of \code{NodeContext} that the network layer and UI/API use for:

\begin{itemize}
  \item pushing blocks into chainstate (\code{add\_block})
  \item answering sync queries (height, hashes, block lookup)
  \item basic peer introspection (who are we connected to?)
\end{itemize}

\begin{listing}[htbp]
\caption{NodeContext block and chain queries}
\label{lst:ch13-block-queries}
\begin{minted}[fontsize=\footnotesize]{rust}
impl NodeContext {
    pub fn get_blockchain(&self)
        -> &BlockchainService {
        &self.blockchain
    }

    pub async fn add_block(
        &self,
        block: &Block
    ) -> Result<()> {
        self.blockchain.add_block(block).await
    }

    pub async fn get_blockchain_height(
        &self
    ) -> Result<usize> {
        self.blockchain.get_best_height().await
    }

    pub async fn get_block_hashes(
        &self
    ) -> Result<Vec<Vec<u8>>> {
        self.blockchain.get_block_hashes().await
    }

    pub async fn get_block(
        &self,
        block_hash: &[u8]
    ) -> Result<Option<Block>> {
        self.blockchain.get_block(block_hash).await
    }

    pub async fn mine_empty_block(
        &self,
        wallet_address: &WalletAddress,
    ) -> Result<Block> {
        miner::mine_empty_block(
            &self.blockchain,
            wallet_address
        ).await
    }

    pub fn get_peer_count(&self) -> Result<usize> {
        use crate::node::GLOBAL_NODES;
        let nodes = GLOBAL_NODES.get_nodes()?;
        Ok(nodes.len())
    }

    pub fn get_peers(&self)
        -> Result<Vec<SocketAddr>> {
        use crate::node::GLOBAL_NODES;
        let nodes = GLOBAL_NODES.get_nodes()?;
        Ok(nodes.into_iter()
            .map(|n| n.get_addr())
            .collect())
    }
}
\end{minted}
\end{listing}

% ──────────────────────────────────────────────────────────────
\section{The Mempool Primitives}
\label{sec:ch13-mempool}

The mempool module is deliberately small: it is a thin wrapper around \code{GLOBAL\_MEMORY\_POOL} plus a UTXO-side flag update so outputs can be marked ``in mempool''.

\begin{listing}[htbp]
\caption{Mempool helpers}
\label{lst:ch13-mempool}
\begin{minted}[fontsize=\footnotesize]{rust}
use crate::error::Result;
use crate::node::GLOBAL_MEMORY_POOL;
use crate::{BlockchainService, Transaction, UTXOSet};

pub async fn add_to_memory_pool(
    tx: Transaction,
    blockchain_service: &BlockchainService,
) -> Result<()> {
    // Store in pool + mark outputs
    // "in mempool" to prevent double-spend
    GLOBAL_MEMORY_POOL.add(tx.clone())?;

    let utxo_set = UTXOSet::new(
        blockchain_service.clone());
    utxo_set.set_global_mem_pool_flag(&tx, true)
        .await?;

    Ok(())
}

pub async fn remove_from_memory_pool(
    tx: Transaction,
    blockchain: &BlockchainService,
) {
    // Remove from pool + clear "in mempool" flags
    let _ = GLOBAL_MEMORY_POOL.remove(tx.clone());
    let utxo_set = UTXOSet::new(
        blockchain.clone());
    let _ = utxo_set
        .set_global_mem_pool_flag(&tx, false)
        .await;
}

pub fn transaction_exists_in_pool(
    tx: &Transaction
) -> bool {
    GLOBAL_MEMORY_POOL.contains_transaction(tx)
        .unwrap_or(false)
}
\end{minted}
\end{listing}

% ──────────────────────────────────────────────────────────────
\section{Mining Pipeline}
\label{sec:ch13-mining}

Mining is triggered by \textbf{policy} (\code{should\_trigger\_mining}) and executed by \code{process\_mine\_block}. The output of mining is a new block which is then announced to peers using \code{INV(block\_hash)}.

The mining trigger $\to$ mining $\to$ relay flow:

\begin{minted}[fontsize=\footnotesize,linenos=false]{text}
NodeContext::submit_transaction_for_mining
  |
  | if should_trigger_mining()
  v
prepare_mining_utxo -> txs + coinbase
  |
  v
process_mine_block -> blockchain.mine_block(txs)
  |
  | remove txs from mempool
  v
broadcast_new_block -> send_inv(op=Block, items=[block_hash])
\end{minted}

\begin{listing}[htbp]
\caption{Mining trigger and block construction}
\label{lst:ch13-mining-trigger}
\begin{minted}[fontsize=\footnotesize]{rust}
const TRANSACTION_THRESHOLD: usize = 3;

pub fn should_trigger_mining() -> bool {
    // Policy: mine if miner configured AND
    // mempool size exceeds threshold
    GLOBAL_MEMORY_POOL.len().unwrap_or(0)
        >= TRANSACTION_THRESHOLD
        && GLOBAL_CONFIG.is_miner()
}

pub async fn prepare_mining_utxo(
    mining_address: &WalletAddress,
) -> Result<Vec<Transaction>> {
    // Snapshot mempool and validate each tx's inputs
    // are still unspent. This prevents mining blocks
    // with already-spent inputs when a competing block
    // has been accepted during the race between miners.
    let txs = GLOBAL_MEMORY_POOL.get_all()?;
    let mut valid_txs = Vec::new();
    for tx in txs {
        if tx.is_coinbase() { continue; }
        // ... (validate each input exists in UTXO tree)
        // Stale transactions are removed from mempool
    }
    valid_txs.push(
        Transaction::new_coinbase_tx(mining_address)?);
    Ok(valid_txs)
}
\end{minted}
\end{listing}

Once we have a validated transaction set, \code{process\_mine\_block} performs the actual proof-of-work, clears the mempool, and announces the result to peers:

\begin{listing}[htbp]
\caption{Mining execution and broadcast}
\label{lst:ch13-mining-exec}
\begin{minted}[fontsize=\footnotesize]{rust}
pub async fn process_mine_block(
    txs: Vec<Transaction>,
    blockchain: &BlockchainService,
) -> Result<Block> {
    // Prevent concurrent mining
    if MINING_IN_PROGRESS.compare_exchange(
        false,
        true,
        Ordering::SeqCst,
        Ordering::SeqCst
    ).is_err() {
        return Err("Mining already in progress".into());
    }
    reset_mining_cancellation();

    let result = async {
        // 1) Mine: PoW + chainstate integration
        // (with stale-mining protection)
        let new_block =
            blockchain.mine_block(&txs).await?;

        // CRITICAL: Do NOT cancel after block creation
        // the block is in our chain and MUST be
        // broadcast, or we create a permanent
        // minority fork.

        // 2) Remove confirmed txs from mempool
        for tx in &txs {
            remove_from_memory_pool(
                tx.clone(),
                blockchain
            ).await;
        }

        // 3) Announce block hash to peers (INV-first)
        broadcast_new_block(&new_block).await?;
        Ok(new_block)
    }.await;

    // Always release lock
    MINING_IN_PROGRESS.store(false, Ordering::SeqCst);
    result
}

pub async fn broadcast_new_block(
    block: &Block
) -> Result<()> {
    let my_addr = GLOBAL_CONFIG.get_node_addr();
    let nodes = GLOBAL_NODES.get_nodes()?;
    for node in nodes {
        if node.get_addr() != my_addr {
            let addr = node.get_addr();
            let hash = block.get_hash_bytes();
            tokio::spawn(async move {
                send_inv(&addr, OpType::Block, &[hash])
                    .await;
            });
        }
    }
    Ok(())
}
\end{minted}
\end{listing}

% ──────────────────────────────────────────────────────────────
\section{Peer Set}
\label{sec:ch13-peer-set}

The peer set is a thread-safe wrapper around a \code{HashSet<Node>} guarded by an \code{RwLock}.

\begin{listing}[htbp]
\caption{Peer data structures}
\label{lst:ch13-peers}
\begin{minted}[fontsize=\footnotesize]{rust}
use crate::error::Result;
use std::collections::HashSet;
use std::net::SocketAddr;
use std::sync::RwLock;

#[derive(Clone, Debug, Eq, PartialEq, Hash)]
pub struct Node { addr: SocketAddr }

impl Node {
    fn new(addr: SocketAddr) -> Self {
        Node { addr }
    }
    pub fn get_addr(&self) -> SocketAddr {
        self.addr
    }
}

// Thread-safe peer set: many reads, occasional writes
pub struct Nodes {
    inner: RwLock<HashSet<Node>>
}

impl Nodes {
    pub fn new() -> Self {
        Nodes {
            inner: RwLock::new(HashSet::new())
        }
    }

    pub fn add_node(&self, addr: SocketAddr)
        -> Result<()> {
        self.inner.write()?
            .insert(Node::new(addr));
        Ok(())
    }

    pub fn add_nodes(
        &self,
        nodes: HashSet<SocketAddr>
    ) -> Result<()> {
        let mut inner = self.inner.write()?;
        for node in nodes {
            inner.insert(Node::new(node));
        }
        Ok(())
    }

    pub fn evict_node(&self, addr: &SocketAddr)
        -> Result<bool> {
        Ok(self.inner.write()?
            .remove(&Node::new(*addr)))
    }

    pub fn get_nodes(&self)
        -> Result<Vec<Node>> {
        Ok(self.inner.read()?
            .iter()
            .cloned()
            .collect())
    }

    pub fn node_is_known(&self, addr: &SocketAddr)
        -> Result<bool> {
        Ok(self.inner.read()?
            .iter()
            .any(|x| x.get_addr() == *addr))
    }
}

impl Default for Nodes {
    fn default() -> Self { Self::new() }
}
\end{minted}
\end{listing}

% ──────────────────────────────────────────────────────────────
\section{Integration Point: Network Routing Uses NodeContext}
\label{sec:ch13-integration}

The network router is responsible for \textbf{decoding} and \textbf{dispatching}, and then calling into \code{NodeContext} for stateful operations.

At the key boundary:

\begin{itemize}
  \item \code{Package::Tx} $\to$ \code{NodeContext::process\_transaction(...)}
  \item \code{Package::Block} $\to$ \code{NodeContext::add\_block(...)}
\end{itemize}

This is the ``runtime wiring'' that makes the node feel like a single system rather than a set of disconnected modules.

% ──────────────────────────────────────────────────────────────
\section{Exercises}
\label{sec:ch13-exercises}

\begin{enumerate}
  \item \textbf{Request Flow Tracing} --- Pick any REST API endpoint (e.g., submit transaction) and trace the request from the HTTP handler through \code{NodeContext} to the final blockchain state change. Identify every module boundary the request crosses.

  \item \textbf{Subsystem Interaction Map} --- Draw a diagram showing how \code{NodeContext} interacts with each subsystem: chain state, mempool, network, mining, and validation. For each interaction, label the function called and the data exchanged.
\end{enumerate}

% ──────────────────────────────────────────────────────────────
\section{Summary: What to Remember}
\label{sec:ch13-summary-remember}

\begin{itemize}
  \item \code{NodeContext} is a façade: it routes high-level requests into mempool, mining, relay, and chainstate.
  \item \code{process\_transaction} is intentionally \textbf{fast}: it accepts to mempool and then spawns relay/mining in the background.
  \item Mining is policy-driven (\code{should\_trigger\_mining}) and ends by relaying an \code{INV} for the new block.
  \item The peer set is a simple \code{RwLock<HashSet<Node>>} and is used by both relay and mining.
  \item Network messages are dispatched to \code{NodeContext} methods, which coordinate subsystems without owning them.
\end{itemize}

% ──────────────────────────────────────────────────────────────
\begin{chaptersummary}
\item We built the \code{NodeContext} that serves as the central coordination point for all blockchain node operations.
\item We unified interactions between blockchain state, transaction mempool, network operations, mining, and validation behind a single API.
\item We saw how the node module abstracts the complexity of coordinating multiple subsystems into clean, testable interfaces.
\item We traced the transaction and block ingestion paths from network packets to persistent blockchain state.
\item We understood how mining policy, peer relay, and consensus validation integrate into a coherent runtime system.
\end{chaptersummary}

In the next chapter, we build the wallet system --- key generation, address derivation, and transaction signing --- that gives users the ability to hold and spend cryptocurrency.
